\subsection{Kỹ năng chuyên môn}

\subsubsection{Kỹ năng lập trình C++}
Dự án game Tetris yêu cầu vận dụng linh hoạt các kiến thức lập trình cơ bản và nâng cao:
\begin{itemize}
    \item \textbf{Cấu trúc dữ liệu và mảng hai chiều:} Ứng dụng ma trận 2D để mô phỏng bảng chơi (grid board) và định nghĩa tọa độ, trạng thái xoay của 7 loại khối gạch Tetromino cơ bản.
    \item \textbf{Xây dựng Game Loop và xử lý thời gian thực:} Cài đặt vòng lặp chính của trò chơi, bắt sự kiện điều khiển phím từ người chơi và điều tiết chu kỳ rơi tự do của khối gạch.
    \item \textbf{Thuật toán kiểm tra va chạm và xóa hàng:}
    \begin{itemize}
        \item Cài đặt kiểm tra va chạm giữa khối gạch đang di chuyển với đường biên và các khối đã cố định trên bảng.
        \item Xây dựng logic quét và xóa các hàng đã lấp đầy, dồn các hàng phía trên xuống, đồng thời giảm thời gian delay (tăng tốc độ rơi) để tăng độ khó khi người chơi ghi điểm.
    \end{itemize}
    \item \textbf{Kỹ thuật gỡ lỗi (Debugging):} Sử dụng công cụ gỡ lỗi của IDE để theo dõi luồng thực thi, kiểm tra biến và bắt các lỗi tràn mảng hoặc sai lệch logic.
\end{itemize}

\subsubsection{Kỹ năng sử dụng Overleaf}
Overleaf là môi trường soạn thảo trực tuyến giúp chuẩn hóa tài liệu báo cáo của dự án:
\begin{itemize}
    \item \textbf{Cộng tác trực tuyến theo thời gian thực:} Cho phép nhiều thành viên cùng đóng góp và chỉnh sửa báo cáo đồng thời mà không bị đè dữ liệu.
    \item \textbf{Biên dịch tài liệu chuẩn học thuật:} Tự động hóa quá trình căn lề, đánh số trang, tạo mục lục và quản lý tài liệu tham khảo chuẩn chỉnh.
    \item \textbf{Cấu trúc báo cáo theo module:} Sử dụng lệnh \texttt{\textbackslash input} để kết nối các tệp \texttt{.tex} thành phần từ mã nguồn vào văn bản hoàn chỉnh một cách khoa học.
\end{itemize}